% University of Helsinki
% Department of Mathematics and Statistics
% Template for Bachelor's or Master's thesis
%
% This template is intended as a starting point for a thesis;
% feel free to customize it or use your own template.
% See the comments for hints and customization.
%
% Initial version: 29 August 2024 / Petri Laarne


% Add the 'twoside, openright' options when producing a proper printed version
\documentclass[a4paper, 11pt]{report}

\usepackage[hidelinks]{hyperref} % Always load this: clickable links and PDF reader navigation bar
\usepackage{lastpage} % Used to determine the total page count
\usepackage{calc} % For layout calculations involving TeX length units (template title page)
\usepackage{graphicx} % For embedding image files (JPEG, PNG, PDF)
\usepackage{subcaption} % For subfigures
\usepackage{tikz} % For drawing diagrams and other complex vector graphics

% If you use the LuaLaTeX or XeLaTeX compiler, comment out the following line.
% These compilers offer better support for Unicode, so you should use them if possible.
\usepackage[T1]{fontenc}

% Pass the language of the thesis as the LAST option in the list below.
% This changes the language used for dates and words like 'Chapter' or 'Bibliography'.
% You can remove the languages that do not appear in your work.
\usepackage[english,finnish]{babel}

\usepackage{csquotes} % This package is suggested by 'babel'

% Bibliography. Change the latter command if you modify the .bib file name.
% \usepackage[backend=biber]{biblatex}
\usepackage{biblatex}
\addbibresource{lähteet.bib}

\usepackage{mathtools} % Extended mathematics support (this package extends amsmath)
\usepackage{amssymb} % More mathematical symbols
\usepackage{amsfonts} % Blackboard bold mathematics font

% Theorem environments
\usepackage{amsthm}
\theoremstyle{plain}
\newtheorem{theorem}{Lause}[section] % Theorems are numbered within chapter and section
\newtheorem{lemma}[theorem]{Lemma} % Lemmas etc. follow the same numbering as theorems
\theoremstyle{definition}
\newtheorem{definition}[theorem]{Määritelmä}
\theoremstyle{remark}
\newtheorem{remark}[theorem]{Huomautus}

% Some examples of custom mathematical commands
\newcommand{\RR}{\mathbb{R}}
\newcommand{\NN}{\mathbb{N}}
\newcommand{\QQ}{\mathbb{Q}}
\newcommand{\ZZ}{\mathbb{Z}}
\newcommand{\dx}{\mathop{}\!\mathrm{d} x} % Nice spacing (and upright d) for dx in integrals

% Finnish notation for evaluating integrals
% Instead of \int_a^b syntax you must use it as \sij{a}{b}
% (To enable the sub-/superscript syntax would require some more complicated TeX wizardry.)
\newcommand{\sij}[2]{\mathop{\bigg/}\nolimits_{\hspace{-0.65em}#1}^{\hspace{0.05em}#2}\hspace{-0.2em}}


\begin{document}

%
% TITLE PAGE
%
\begin{titlepage}

\noindent
\begin{minipage}[b]{\textwidth-45pt}
% University of Helsinki\\
% Faculty of Science\\
Matemaattisten tieteiden kandiohjelma\\
Helsingin yliopisto\\
Matemaattis-luonnontieteellinen tiedekunta\\
\end{minipage}
\hfill\includegraphics[width=45pt]{HYlogo.pdf}

\vspace{4pt}\hrule\vfill

\begin{center}
% Bachelor's thesis\\[8pt]
Kandidaatintutkielma\\[8pt]

{\huge\bfseries Mitta jotainjotain}\\[8pt]

Aaro Koivukangas
\end{center}

\vfill\hrule\vspace{4pt}

\noindent
Ohjaaja:\hfill Palautuskuukausi ja -vuosi\\
Petri Ola
\end{titlepage}

%
% ABSTRACT
%
% This command starts a new page.
% In a print version ('twoside' option passed to \documentclass), the page is on the right side
\pagenumbering{gobble}
\cleardoublepage

\noindent
\textbf{Otsikko:} Mitta jotainjotain\\ % Otsikko / Titel
\textbf{Tekijä:} Aaro Koivukangas\\ % Tekijä / Författare
\textbf{Suorituskuukausi ja -vuosi} Palautuskuukausi ja -vuosi\\ % Aika / Datum
\textbf{Sivumäärä:}~\pageref*{LastPage}~s.\\[1em] % Sivumäärä: XX~s. / Sidoantal: XX~s.
% The page count is updated automatically, but it might require two compilations to be correct.
% If you customize the backmatter, this might become incorrect - check and correct manually if needed.

% If you include an abstract in a different language, wrap it in a foreignlanguage block
% \begin{foreignlanguage}{english}
\begin{foreignlanguage}{finnish}
\noindent\textbf{Tiivistelmä:}\\ % Abstract

Loorem ipsum ja sitä rataa. Kanditutkielman teko-ohje:

\url{https://studies.helsinki.fi/ohjeet/artikkeli/tutkielma-ja-kypsyysnayte-kandiohjelmassa?degree_programme_code=KH50_001}

\end{foreignlanguage}


\vfill

\noindent\textbf{Avainsanat:} Mitta, Integraali % Keywords




%
% TABLE OF CONTENTS
%
\cleardoublepage
\tableofcontents


%
% MAIN TEXT STARTS HERE
%

% The \pagenumbering{arabic} command must be immediately after the first chapter title.
% Do not repeat it for other chapters.
\chapter{Introduction}
\pagenumbering{arabic}

\section{Compiling the document}

\section{Mathematics}

\section{Theorems, definitions, and proofs}


\begin{definition}
A triple $(a, b, c)$ of positive integers is called \emph{Pythagorean} if $a^2 + b^2 = c^2$.
\end{definition}

It is possible to give a name to such a block.

\begin{lemma}[A very silly lemma]
There exists at least one Pythagorean triple.
\end{lemma}
\begin{proof}
Pick $(a, b, c) = (3, 4, 5)$.
\end{proof}

If you want the block name to refer to a result in a published work,
you need to wrap the citation command in \verb|{{}}|.
This is due to some unfortunate \TeX nical details.

\begin{theorem}[{{\cite[Theorem~0.5]{wiles}}}]
There does not exist a triple $(a, b, c)$ of positive integers such that $a^3 + b^3 = c^3$.
\end{theorem}
\begin{proof}
Left as an exercise to the reader.
\end{proof}

\begin{remark}
The last proof was a joke.
\end{remark}



\section{Citing literature}

\nocite{*}
Jeejee

%
% BIBLIOGRAPHY
%
% The [heading=bibintoc] option adds the bibliography to table of contents
\printbibliography[heading=bibintoc]




%
% APPENDICES (if necessary)
%
\appendix

\chapter{Jottainnii lisätietoja}

\end{document}
